\documentclass[11pt,a4paper,oneside]{article}\usepackage[]{graphicx}\usepackage[]{color}
%% maxwidth is the original width if it is less than linewidth
%% otherwise use linewidth (to make sure the graphics do not exceed the margin)
\makeatletter
\def\maxwidth{ %
  \ifdim\Gin@nat@width>\linewidth
    \linewidth
  \else
    \Gin@nat@width
  \fi
}
\makeatother

\definecolor{fgcolor}{rgb}{0.345, 0.345, 0.345}
\newcommand{\hlnum}[1]{\textcolor[rgb]{0.686,0.059,0.569}{#1}}%
\newcommand{\hlstr}[1]{\textcolor[rgb]{0.192,0.494,0.8}{#1}}%
\newcommand{\hlcom}[1]{\textcolor[rgb]{0.678,0.584,0.686}{\textit{#1}}}%
\newcommand{\hlopt}[1]{\textcolor[rgb]{0,0,0}{#1}}%
\newcommand{\hlstd}[1]{\textcolor[rgb]{0.345,0.345,0.345}{#1}}%
\newcommand{\hlkwa}[1]{\textcolor[rgb]{0.161,0.373,0.58}{\textbf{#1}}}%
\newcommand{\hlkwb}[1]{\textcolor[rgb]{0.69,0.353,0.396}{#1}}%
\newcommand{\hlkwc}[1]{\textcolor[rgb]{0.333,0.667,0.333}{#1}}%
\newcommand{\hlkwd}[1]{\textcolor[rgb]{0.737,0.353,0.396}{\textbf{#1}}}%

\usepackage{framed}
\makeatletter
\newenvironment{kframe}{%
 \def\at@end@of@kframe{}%
 \ifinner\ifhmode%
  \def\at@end@of@kframe{\end{minipage}}%
  \begin{minipage}{\columnwidth}%
 \fi\fi%
 \def\FrameCommand##1{\hskip\@totalleftmargin \hskip-\fboxsep
 \colorbox{shadecolor}{##1}\hskip-\fboxsep
     % There is no \\@totalrightmargin, so:
     \hskip-\linewidth \hskip-\@totalleftmargin \hskip\columnwidth}%
 \MakeFramed {\advance\hsize-\width
   \@totalleftmargin\z@ \linewidth\hsize
   \@setminipage}}%
 {\par\unskip\endMakeFramed%
 \at@end@of@kframe}
\makeatother

\definecolor{shadecolor}{rgb}{.97, .97, .97}
\definecolor{messagecolor}{rgb}{0, 0, 0}
\definecolor{warningcolor}{rgb}{1, 0, 1}
\definecolor{errorcolor}{rgb}{1, 0, 0}
\newenvironment{knitrout}{}{} % an empty environment to be redefined in TeX

\usepackage{alltt}
\usepackage{amsmath,amsthm,amsfonts,amssymb}
\usepackage{pst-eucl,pstricks,pstricks-add}
\usepackage[utf8]{inputenc}
%\usepackage[latin1]{inputenc}
\usepackage[spanish,activeacute]{babel}
\usepackage[a4paper,margin=2cm]{geometry}
\usepackage{times}
\usepackage[T1]{fontenc}
\usepackage{titlesec}
\usepackage{color}
\usepackage{url}
\usepackage{float}
\usepackage{cite}
\usepackage{graphicx}
\usepackage{multicol}
\usepackage{lmodern}
\usepackage{setspace}
%\doublespace %para doble espacio
\onehalfspace %para espacio y medio
\newcommand{\code}[1]{\fcolorbox{blue!80}{gray!10}{#1}}
\parindent=0mm
\IfFileExists{upquote.sty}{\usepackage{upquote}}{}
\begin{document}
\begin{minipage}[d]{115mm}
\includegraphics[height=2.5cm, width=7.5cm]{logo_ult.pdf}
\end{minipage}
\begin{minipage}[d]{70mm}
\textsf{\textbf{\sc \tiny {\bf QUITO}: Fernando Oviedo E8-65 y José Barba}}\\
%\textsf{\textbf{\sc \tiny \qquad  Sector El Dorado}}\\
\textsf{\textbf{\sc \tiny {\bf CONTACTOS}: (02) 2 559-703 / (02) 2 602-560\\ Movistar: 0998428362 / 0998890021}}
\end{minipage}\\
\vspace{0.4cm}

\begin{center}
{\bf \Large Curso: R Nivel Básico}
\end{center}

\section{Descripción}
  
El curso de R Nivel Básico tiene como objetivo principal proporcionar al estudiante una visión general del entorno del lenguaje de programación R y sus aplicaciones. Al final del curso, el estudiante se encontrará en la capacidad de crear sus propias funciones, importar y exportar datos.

\section{Contenidos del curso}

\begin{multicols}{2}

\begin{enumerate}
  \item {\Large Introducción}
  \begin{enumerate}
    \item[1.1] Funcionamiento
    \item[1.2] Ventajas $\&$ Desventajas
  \end{enumerate}
  \item {\Large Instalación y actualización}
  \begin{enumerate}
    \item[2.1] Entorno de trabajo
    \item[2.2] Instalación de paquetes
    \item[2.3] Actualización de paquetes
    \item[2.4] Uso de Git $\&$ Github
    \item[2.5] Instalación RRO
  \end{enumerate}
  \item {\Large Interfaces de usuario}
  \begin{enumerate}
    \item[3.1] RStudio
    \item[3.2] R Analytic Flow
  \end{enumerate}
  \item {\Large Estructura de datos}
    \begin{enumerate}
    \item[4.1] Vectores atómicos $\&$ Listas
    \item[4.2] Matrices $\&$ Arrays
    \item[4.3] Factores
    \item[4.4] Data Frames $\&$ Data Tables
    \item[4.5] Coerción
  \end{enumerate}
  \item {\Large Lectura y escritura de datos}
  \begin{enumerate}
    \item[5.1] Ingreso mediante el teclado
    \item[5.2] Lectura de archivos
    \item[5.3] Escritura de archivos
    \item[5.4] Accediendo a datos de la web
    \item[5.6] Paquetes \texttt{haven}, \texttt{readr}, \texttt{readxl}
  \end{enumerate}
  \item {\Large Funciones}
    \begin{enumerate}
    \item[6.1] Funciones de R
    \item[6.2] Estructuras de control
    \item[6.3] Creación de funciones
    \item[6.4] Ámbitos léxicos (Lexical Scoping)
    \item[6.5] Ambientes (Environments)
  \end{enumerate}
\end{enumerate}

\end{multicols}

\section{Duración del curso}
  
El curso tiene una duración de 20 horas, mismas que serán distribuidas en 4 sesiones (Sábados, a partir del 07 de Noviembre) en el horario de 8:00 a 13:00, en el Laboratorio 04 del Departamento de Matemática.

\section{Costos}

El costo del curso se detalla a continuación:

\begin{center}
\begin{tabular}{|c|c|c|}\hline
{\bf Grupo} & {\bf Costo}\\ \hline
Estudiantes EPN & 40 dólares \\ \hline
Particulares & 50 dólares\\ \hline
\end{tabular}
\end{center}

\section{Lugar}

Escuela Politécnica Nacional, Departamento de Matemática, Edificio Administrativo, Séptimo Piso, Laboratorio 04.

\end{document}
